\section{From raw experiments to a measurable predictive quotient}\label{sec:quotient}
We give a sufficient realization result at the raw level. It does not assert that every standard Borel predictive equivalence relation has a compact or even standard Borel quotient.

For each phase let $X_t$ be compact metrizable. Let
$K_t^a(x,dx'\,dy)$ be a positive normalized Borel physical kernel with standard Borel reports. Fix on each report space a compatible Polish topology and a countable bounded continuous measure-determining family containing constants and closed under rational linear combinations and products. Let $\mathcal V_t\subset C(X_t)$ be a separable closed linear subspace containing constants and the conditional expectations of a determining family of executable future tests. Assume it contains every pullback
\begin{equation}\label{eq:raw_pullback}
 x\longmapsto\int \varphi(y)f(x')K_t^a(x,dx'\,dy),
 \qquad f\in\mathcal V_{t+1},
\end{equation}
for legal $a$ and each function $\varphi$ in the chosen report family. In particular these pullbacks are continuous. Assume the declared future-test system is determined by these tests and their successive report-contingent continuations; equality on them must not discard a legal action or an observable resource coordinate. Phase-dependent legality is retained separately. The terminal task has continuous affine expectation in the resulting coordinates. These are verifiable closure conditions on the raw experiment, not assumptions of a Fisher metric or an available exact posterior register.

Choose a countable norm-dense family $f_{t,i}$ in the unit ball of $\mathcal V_t$, including constants, and define
\begin{equation}\label{eq:raw_quotient}
 Q_t:\mathcal P(X_t)\longrightarrow[-1,1]^{\mathbb N},\qquad
 Q_t(p)=(\int f_{t,i}\,dp)_{i\ge1},\qquad
 \mathsf K_t=Q_t(\mathcal P(X_t)).
\end{equation}
The dense family may be chosen from the rational span of determining executable tests; equality of its coordinates is therefore equality of their predictions, not an extra observable added to the experiment.

\begin{theorem}[Predictive realization and minimality]\label{thm:quotient}
Under the preceding raw hypotheses, $\mathsf K_t$ is compact metrizable and convex. Two physical distributions have the same $Q_t$ precisely when their declared executable future tests agree. The actual posterior of every legal finite-history strategy projects to a Borel $\mathsf K_t$-valued predictive state. Report and update kernels on $\mathsf K_t$ admit jointly Borel versions, independent of the representing distribution up to the relevant report-null sets, and satisfy \eqref{eq:affine_instrument}. Conditional mixing is realized by barycenters.

Every measurable statistic determining all these future tests factors the projected predictive state through a Borel map; that factor is unique on the statistic's attained values. Conversely, the projected state determines the tests. This is minimality for the fixed preparation/test interface, not simultaneous frequentist sufficiency for all inaccessible parameters. If the induced acquired-law kernels are weakly continuous, all the hypotheses concerning predictive realization in Theorem~\ref{thm:completion} hold.
\end{theorem}
\begin{proof}
Weak compactness of $\mathcal P(X_t)$ and continuity of the coordinates make $\mathsf K_t$ compact in the countable product; it is convex by linearity. Equality of dense coordinates gives equality on $\mathcal V_t$. The report pullbacks with $f=1$ determine the report law. With general $f$ they determine the finite signed measures
$B\mapsto\int_{X_t\times X_{t+1}\times B}f(x')p(dx)K_t^a(x,dx'\,dy)$.
The measure-determining report family and a monotone-class argument extend the identities to all Borel report sets. After disintegration, equal incoming coordinates therefore give equal outgoing conditional coordinates outside one report-null set, chosen common for the countable family. Induction over the finite continuation proves equality of the declared future tests, including report-contingent choices. The reverse implication holds because the determining tests are executable. Continuity of arbitrary report-contingent test expectations is not needed for this implication.

For measurability, a continuous map from a compact metric space onto a compact metric image has a Borel section here. An explicit construction is useful. Cover $\mathcal P(X_t)$ at level $m$ by finitely many closed balls of radius $2^{-m}$. Select the least ball whose intersection with the previously selected compact intersections still meets the fiber $Q_t^{-1}(q)$. The condition is Borel in $q$, since the image of each such compact intersection is closed. The nested nonempty compact intersections have diameters tending to zero, and their unique limiting point is a Borel selection $p(q)$ in the fiber.

Apply a jointly measurable conditional-probability construction to the raw kernel with input $(p(q),a)$. Such a construction follows, for example, by conditioning on successive finite partitions generating the report sigma field: conditional expectations of a countable dense family in $C(X_{t+1})$ converge by the martingale theorem, determine a probability measure on the compact state space almost everywhere, and have measurable limits. Assign an arbitrary fixed probability where the convergence or consistency tests fail. Project this posterior by $Q_{t+1}$ to define $F_t(q,a,y)$, and use the raw report marginal to define $J_t(q,a)$. The preceding signed-measure identities prove representative independence almost everywhere. The initial and all later history spaces are standard Borel, so this construction also gives Borel projected posteriors along actual histories.

For a probability $\nu$ on $\mathsf K_t$, mix its selected representatives: $p=\int p(q)\nu(dq)$. Then $Q_t(p)=\bar\nu$. Linearity of the unnormalized raw kernel in $p$, followed by the signed-measure identities, gives both formulas \eqref{eq:affine_instrument}. Equality on coordinate functions is enough: it identifies each conditional barycenter in $\mathsf K_{t+1}$ and hence every continuous affine function. This is mathematical realizability as a conditional physical mixture. It does not make re-preparation free or assert that every point of the compact carrier is reached under the chosen exploration. In a feasible allocation flow the induction of Lemma~\ref{lem:calibration} supplies the actual conditional mixtures under the deployed machine.

Finally, let a statistic $T(h)$ determine the tests, with Borel coordinate maps $b_i$ such that $Q_t(p_h)_i=b_i(T(h))$. The vector $b=(b_i)$ is Borel. Since $\mathsf K_t$ is closed in the product, replace $b$ by a fixed point outside the Borel set $b^{-1}(\mathsf K_t)$. This defines a Borel factor through $T$ agreeing with the projected state on every attained history where the conditional versions are specified. Its values there are forced by the separating coordinates. No claim that the image of an arbitrary history map is Borel is required. The remaining weak-continuity condition is additional; it is not a consequence of measurability alone.
\end{proof}
In the finite hidden-state model, the test span identifies the simplex of distributions. In the noncommuting instrument model, all future-test expectations are affine in the density matrix and the spanning effects identify that matrix; different classical decompositions of it are correctly identified. For capture of a compactly supported variable, bounded continuous report tests identify its probability law, and the carrier can be infinite-dimensional. These identifications explain why the same abstract theorem covers the three settings without assigning an ambient Fisher or covariance geometry at the outset.
